\begin{abstract}
\noindent 再生数估计误差常被归因于样本量。本文区分当期瞬时再生数 $R_t$ 与跨期均值 $\Rbar$。负二项传播与二项报告下，前者误差随预计报告规模 $m$ 按 $1/m$ 下降；后者含不随单期规模下降的环境方差 $v_R$ 与报告异质性方差 $v_H$。这里的“不可约”针对扩大单期规模而言，并不意味着增加独立观测期数无法降低该误差。由此导出方差平衡规模 $\mx=A/v_R$，将监测设计归结为规模、期数与口径三个杠杆。基于美国 51 个辖区 2020--2026 年流感住院 10,303 个有效州周，流行初期误差关于 $1/D$ 的截距为 0.0991（95\% 区间 0.0837--0.1165），斜率为 2.243（1.494--3.007），经验比值 $b/a=22.6$（13.5--34.9）。实证观察到与理论误差分解相容的收益递减模式，不过尚不足以区分严格非零平台与缓慢衰减，也不能将截距唯一归因于环境方差。报告调整后截距变化为 $-0.0004$，填报比例代理能解释的报告成分有限（不足 1\%）。结果表明跨期趋势估计面临单期误差下限，预警规模应按检验功效单独确定。

\vspace{4pt}
\noindent\textbf{关键词：}再生数；环境随机性；报告异质性；监测设计；方差分解；流感住院
\end{abstract}
